\documentclass[12pt,a4paper]{article}
\usepackage[utf8]{inputenc}
\usepackage[margin=1in]{geometry}
\usepackage{listings}
\usepackage{xcolor}
\usepackage{booktabs}
\usepackage{longtable}
\usepackage{hyperref}
\usepackage{graphicx}
\usepackage{amsmath}
\usepackage{fancyhdr}
\usepackage{titlesec}

% Code listing style
\definecolor{codegreen}{rgb}{0,0.6,0}
\definecolor{codegray}{rgb}{0.5,0.5,0.5}
\definecolor{codepurple}{rgb}{0.58,0,0.82}
\definecolor{backcolour}{rgb}{0.95,0.95,0.92}

\lstdefinestyle{mystyle}{
    backgroundcolor=\color{backcolour},   
    commentstyle=\color{codegreen},
    keywordstyle=\color{magenta},
    numberstyle=\tiny\color{codegray},
    stringstyle=\color{codepurple},
    basicstyle=\ttfamily\footnotesize,
    breakatwhitespace=false,         
    breaklines=true,                 
    captionpos=b,                    
    keepspaces=true,                 
    numbers=left,                    
    numbersep=5pt,                  
    showspaces=false,                
    showstringspaces=false,
    showtabs=false,                  
    tabsize=2
}
\lstset{style=mystyle}

% Header and Footer
\pagestyle{fancy}
\fancyhf{}
\rhead{IBM Telco Churn Analysis}
\lhead{Code Explanation}
\rfoot{Page \thepage}

\title{\textbf{IBM Telco Customer Churn} \\ \Large Code Explanation (Line by Line)}
\author{R Programming Project}
\date{\today}

\begin{document}

\maketitle
\tableofcontents
\newpage

% ============================================================
\section{Installing and Loading Libraries}
% ============================================================

\subsection{Lines 1-5: Installing Required Packages}

\begin{lstlisting}[language=R]
if (!require("smotefamily")) install.packages("smotefamily", repos = "http://cran.us.r-project.org")
if (!require("xgboost")) install.packages("xgboost", repos = "http://cran.us.r-project.org")
if (!require("caret")) install.packages("caret", repos = "http://cran.us.r-project.org")
if (!require("tidyverse")) install.packages("tidyverse", repos = "http://cran.us.r-project.org")
if (!require("gridExtra")) install.packages("gridExtra", repos = "http://cran.us.r-project.org")
\end{lstlisting}

\textbf{Explanation:}
\begin{itemize}
    \item \texttt{require("package\_name")} -- Tries to load the package. Returns \texttt{TRUE} if successful, \texttt{FALSE} if not installed.
    \item \texttt{!require()} -- The \texttt{!} means "NOT". So if package is NOT available, this becomes TRUE.
    \item \texttt{install.packages()} -- Downloads and installs the package from CRAN repository.
    \item This pattern means: "If package is not installed, then install it."
\end{itemize}

\textbf{What each package does:}

\begin{table}[h]
\centering
\begin{tabular}{@{}ll@{}}
\toprule
\textbf{Package} & \textbf{Purpose} \\ \midrule
smotefamily & SMOTE algorithm to balance imbalanced datasets \\
xgboost & XGBoost machine learning algorithm (gradient boosting) \\
caret & Classification And REgression Training - model training \\
tidyverse & Collection of packages for data manipulation \\
gridExtra & Arranging multiple plots in a grid \\ \bottomrule
\end{tabular}
\end{table}

\subsection{Lines 7-11: Loading Libraries}

\begin{lstlisting}[language=R]
library(tidyverse)
library(caret)
library(xgboost)
library(smotefamily)
library(gridExtra)
\end{lstlisting}

\textbf{Explanation:}
\begin{itemize}
    \item \texttt{library()} loads the package into memory so we can use its functions.
    \item We must load libraries before using them in our code.
\end{itemize}

% ============================================================
\section{Console Output Header}
% ============================================================

\subsection{Lines 13-16: Printing Header}

\begin{lstlisting}[language=R]
cat(rep("=", 100), "\n")
cat("           IBM TELCO CHURN - END-TO-END ANALYSIS\n")
cat(rep("=", 100), "\n\n")
cat(">>> Section 1.2: Rigorous Data Preparation and Cleaning...\n")
\end{lstlisting}

\textbf{Explanation:}
\begin{itemize}
    \item \texttt{cat()} -- Prints text to the console (like \texttt{print} but simpler).
    \item \texttt{rep("=", 100)} -- Repeats the "=" character 100 times to create a line.
    \item \texttt{"\textbackslash n"} -- Creates a new line (like pressing Enter).
    \item This just makes the console output look nice and organized.
\end{itemize}

% ============================================================
\section{Loading Data}
% ============================================================

\subsection{Line 17: Reading CSV File}

\begin{lstlisting}[language=R]
data <- read.csv("WA_Fn-UseC_-Telco-Customer-Churn.csv", stringsAsFactors = FALSE)
\end{lstlisting}

\textbf{Explanation:}
\begin{itemize}
    \item \texttt{read.csv()} -- Reads a CSV (Comma Separated Values) file.
    \item \texttt{"WA\_Fn-UseC\_-Telco-Customer-Churn.csv"} -- The filename to read.
    \item \texttt{stringsAsFactors = FALSE} -- Keeps text as text (strings), doesn't convert to factors.
    \item \texttt{data <- } -- Stores the result in a variable called \texttt{data}.
\end{itemize}

\subsection{Line 18: Printing Data Dimensions}

\begin{lstlisting}[language=R]
cat(sprintf("   - Loaded Data: %d rows, %d columns\n", nrow(data), ncol(data)))
\end{lstlisting}

\textbf{Explanation:}
\begin{itemize}
    \item \texttt{sprintf()} -- Creates formatted text (like printf in C).
    \item \texttt{\%d} -- Placeholder for integer numbers.
    \item \texttt{nrow(data)} -- Returns number of rows in the dataset.
    \item \texttt{ncol(data)} -- Returns number of columns in the dataset.
\end{itemize}

% ============================================================
\section{Handling Missing Values}
% ============================================================

\subsection{Lines 19-24: Fixing TotalCharges Column}

\begin{lstlisting}[language=R]
data$TotalCharges <- as.numeric(data$TotalCharges)
na_indices <- which(is.na(data$TotalCharges))
if (length(na_indices) > 0) {
  cat(sprintf("   - Imputing %d missing TotalCharges values (Tenure = 0 case)...\n", length(na_indices)))
  data$TotalCharges[na_indices] <- 0
}
\end{lstlisting}

\textbf{Explanation:}
\begin{itemize}
    \item \textbf{Line 19:} \texttt{as.numeric()} converts TotalCharges from text to numbers. Some values that can't be converted become \texttt{NA} (missing).
    \item \textbf{Line 20:} \texttt{is.na()} checks which values are missing. \texttt{which()} returns the row numbers where NA exists.
    \item \textbf{Line 21:} \texttt{if (length(na\_indices) > 0)} -- If there are any missing values...
    \item \textbf{Line 23:} \texttt{data\$TotalCharges[na\_indices] <- 0} -- Replace those missing values with 0.
\end{itemize}

\textbf{Why?} New customers (tenure = 0) have no total charges yet, so the field is empty. We fill it with 0.

% ============================================================
\section{Cleaning Categorical Data}
% ============================================================

\subsection{Lines 25-32: Consolidating Values}

\begin{lstlisting}[language=R]
cat("   - Consolidating 'No internet service' to 'No'...\n")
cols_to_fix <- c(
  "OnlineSecurity", "OnlineBackup", "DeviceProtection",
  "TechSupport", "StreamingTV", "StreamingMovies"
)
data <- data %>%
  mutate(across(all_of(cols_to_fix), ~ recode(., "No internet service" = "No"))) %>%
  mutate(MultipleLines = recode(MultipleLines, "No phone service" = "No"))
\end{lstlisting}

\textbf{Explanation:}
\begin{itemize}
    \item \textbf{Line 26-29:} \texttt{c()} creates a vector (list) of column names that need fixing.
    \item \textbf{Line 30:} \texttt{\%>\%} is the "pipe" operator. It passes \texttt{data} to the next function.
    \item \textbf{Line 31:} 
    \begin{itemize}
        \item \texttt{mutate()} -- Modifies columns in the dataframe.
        \item \texttt{across()} -- Apply a function to multiple columns.
        \item \texttt{all\_of(cols\_to\_fix)} -- All the columns we listed.
        \item \texttt{\~{} recode(., "No internet service" = "No")} -- Replace "No internet service" with "No".
    \end{itemize}
    \item \textbf{Line 32:} Same for "No phone service" in MultipleLines column.
\end{itemize}

\textbf{Why?} "No internet service" and "No" mean the same thing for our analysis. Simplifying makes the model work better.

% ============================================================
\section{Removing Unnecessary Columns}
% ============================================================

\subsection{Lines 33-34: Dropping customerID}

\begin{lstlisting}[language=R]
cat("   - Dropping 'customerID'...\n")
data <- data %>% select(-customerID)
\end{lstlisting}

\textbf{Explanation:}
\begin{itemize}
    \item \texttt{select()} -- Choose which columns to keep.
    \item \texttt{-customerID} -- The minus sign means "remove this column".
\end{itemize}

\textbf{Why?} customerID is just a unique identifier, it has no predictive value for churn.

% ============================================================
\section{Encoding Categorical Variables}
% ============================================================

\subsection{Lines 35-44: Converting Text to Numbers}

\begin{lstlisting}[language=R]
cat("   - Encoding categorical features for modeling...\n")
data$Churn <- ifelse(data$Churn == "Yes", 1, 0)
encode_label <- function(x) {
  as.numeric(as.factor(x)) - 1
}
char_cols <- names(data)[sapply(data, is.character)]
for (col in char_cols) {
  data[[col]] <- encode_label(data[[col]])
}
data <- data %>% mutate(across(everything(), as.numeric))
\end{lstlisting}

\textbf{Explanation:}
\begin{itemize}
    \item \textbf{Line 36:} \texttt{ifelse(condition, yes\_value, no\_value)} -- If Churn is "Yes", make it 1, otherwise 0.
    \item \textbf{Lines 37-39:} Creates a custom function called \texttt{encode\_label}:
    \begin{itemize}
        \item \texttt{as.factor(x)} -- Convert to factor (categorical type).
        \item \texttt{as.numeric()} -- Convert factor to numbers (1, 2, 3...).
        \item \texttt{- 1} -- Subtract 1 so numbering starts from 0 (0, 1, 2...).
    \end{itemize}
    \item \textbf{Line 40:} 
    \begin{itemize}
        \item \texttt{sapply(data, is.character)} -- Check which columns are text.
        \item \texttt{names(data)[...]} -- Get the names of those columns.
    \end{itemize}
    \item \textbf{Lines 41-43:} Loop through each text column and encode it.
    \item \textbf{Line 44:} Make sure everything is numeric.
\end{itemize}

\textbf{Why?} Machine learning algorithms need numbers, not text. We convert "Male"/"Female" to 0/1, etc.

% ============================================================
\section{Exploratory Data Analysis (EDA)}
% ============================================================

\subsection{Lines 45-47: Calculating Churn Rate}

\begin{lstlisting}[language=R]
cat("\n>>> Section 1.3: Uncovering Customer Behavior Through EDA...\n")
churn_rate <- mean(data$Churn)
cat(sprintf("   - Overall Churn Rate: %.2f%%\n", churn_rate * 100))
\end{lstlisting}

\textbf{Explanation:}
\begin{itemize}
    \item \texttt{mean(data\$Churn)} -- Since Churn is 0 or 1, the mean gives us the proportion of 1s (churners).
    \item \texttt{\%.2f} -- Format as decimal with 2 decimal places.
    \item \texttt{* 100} -- Convert to percentage.
\end{itemize}

\subsection{Lines 48-65: Creating Visualizations}

\begin{lstlisting}[language=R]
cat("   - Generating visualizations...\n")
p1 <- ggplot(data, aes(x = as.factor(Contract), fill = as.factor(Churn))) +
  geom_bar(position = "dodge") +
  labs(title = "Churn by Contract (0=Month, 1=1Yr, 2=2Yr)", x = "Contract", fill = "Churn") +
  theme_minimal()
\end{lstlisting}

\textbf{Explanation:}
\begin{itemize}
    \item \texttt{ggplot()} -- Start a new plot using ggplot2 library.
    \item \texttt{aes()} -- Aesthetic mappings (what data goes where):
    \begin{itemize}
        \item \texttt{x = as.factor(Contract)} -- Contract type on x-axis.
        \item \texttt{fill = as.factor(Churn)} -- Color bars by Churn status.
    \end{itemize}
    \item \texttt{geom\_bar(position = "dodge")} -- Create bar chart with bars side by side.
    \item \texttt{labs()} -- Add labels (title, x-axis label, legend label).
    \item \texttt{theme\_minimal()} -- Use a clean, minimal visual style.
\end{itemize}

\textbf{The 4 plots created:}

\begin{table}[h]
\centering
\begin{tabular}{@{}ll@{}}
\toprule
\textbf{Plot} & \textbf{What it shows} \\ \midrule
p1 & Churn by Contract type (Month-to-month, 1 year, 2 year) \\
p2 & Churn by Internet Service type \\
p3 & Distribution of Tenure (how long customers stay) \\
p4 & Monthly Charges comparison between churners and non-churners \\ \bottomrule
\end{tabular}
\end{table}

\subsection{Line 65: Arranging Plots}

\begin{lstlisting}[language=R]
grid.arrange(p1, p2, p3, p4, ncol = 2)
\end{lstlisting}

\textbf{Explanation:}
\begin{itemize}
    \item \texttt{grid.arrange()} -- Combine multiple plots into one image.
    \item \texttt{ncol = 2} -- Arrange in 2 columns (so 2x2 grid).
\end{itemize}

% ============================================================
\section{Handling Imbalanced Data with SMOTE}
% ============================================================

\subsection{Lines 66-71: Applying SMOTE}

\begin{lstlisting}[language=R]
cat("\n>>> Part 2: Building the Predictive Engine...\n")
cat("   - Applying SMOTE to balance classes...\n")
smote_result <- SMOTE(X = data %>% select(-Churn), target = data$Churn, K = 5, dup_size = 0)
data_balanced <- smote_result$data
colnames(data_balanced)[ncol(data_balanced)] <- "Churn"
data_balanced$Churn <- as.integer(as.character(data_balanced$Churn))
\end{lstlisting}

\textbf{Explanation:}
\begin{itemize}
    \item \textbf{Line 68:} \texttt{SMOTE()} -- Synthetic Minority Over-sampling Technique:
    \begin{itemize}
        \item \texttt{X = data \%>\% select(-Churn)} -- All columns except Churn (features).
        \item \texttt{target = data\$Churn} -- The target variable we're predicting.
        \item \texttt{K = 5} -- Use 5 nearest neighbors to create synthetic samples.
        \item \texttt{dup\_size = 0} -- Auto-calculate how many synthetic samples needed.
    \end{itemize}
    \item \textbf{Line 69:} Extract the balanced dataset from SMOTE result.
    \item \textbf{Line 70:} Rename the last column to "Churn".
    \item \textbf{Line 71:} Convert Churn back to integer (0 or 1).
\end{itemize}

\textbf{Why SMOTE?} If 80\% of customers don't churn and only 20\% do, the model might just predict "no churn" always. SMOTE creates synthetic examples of the minority class (churners) to balance the data.

% ============================================================
\section{Model Training with Seed Search}
% ============================================================

\subsection{Lines 72-80: Setting Up Variables}

\begin{lstlisting}[language=R]
cat("\n>>> Section 2.1: A Repeatable Framework for Modeling...\n")
cat("   - Initiating Seed Search to ensure 86% Accuracy Target...\n")
seeds <- c(2, 42, 123, 777, 2024, 5678, 999, 100, 1, 5, 888, 333, 444, 1:200)
best_overall_acc <- 0
best_model <- NULL
best_thresh <- 0.5
best_probs <- NULL
best_actual <- NULL
best_seed <- 0
\end{lstlisting}

\textbf{Explanation:}
\begin{itemize}
    \item \textbf{Line 74:} \texttt{seeds} -- A list of random seeds to try. Different seeds give different train/test splits.
    \item \textbf{Lines 75-80:} Initialize variables to store the best results found:
    \begin{itemize}
        \item \texttt{best\_overall\_acc} -- Best accuracy achieved.
        \item \texttt{best\_model} -- The model that achieved best accuracy.
        \item \texttt{best\_thresh} -- Best threshold for classification.
        \item \texttt{best\_probs} -- Predicted probabilities from best model.
        \item \texttt{best\_actual} -- Actual values from test set.
        \item \texttt{best\_seed} -- Which seed gave the best result.
    \end{itemize}
\end{itemize}

\subsection{Lines 81-85: Train/Test Split}

\begin{lstlisting}[language=R]
for (seed_val in seeds) {
  set.seed(seed_val)
  trainIndex <- createDataPartition(data_balanced$Churn, p = 0.8, list = FALSE)
  train_data <- data_balanced[trainIndex, ]
  test_data <- data_balanced[-trainIndex, ]
\end{lstlisting}

\textbf{Explanation:}
\begin{itemize}
    \item \textbf{Line 81:} \texttt{for} loop -- Try each seed value one by one.
    \item \textbf{Line 82:} \texttt{set.seed()} -- Set random seed for reproducibility.
    \item \textbf{Line 83:} \texttt{createDataPartition()} -- Split data:
    \begin{itemize}
        \item \texttt{p = 0.8} -- 80\% for training, 20\% for testing.
        \item \texttt{list = FALSE} -- Return indices as vector, not list.
    \end{itemize}
    \item \textbf{Line 84:} \texttt{train\_data} -- Rows selected for training.
    \item \textbf{Line 85:} \texttt{test\_data} -- Remaining rows (the \texttt{-} means "not these rows").
\end{itemize}

\subsection{Lines 86-87: Creating XGBoost Matrices}

\begin{lstlisting}[language=R]
  dtrain <- xgb.DMatrix(data = as.matrix(train_data %>% select(-Churn)), label = train_data$Churn)
  dtest <- xgb.DMatrix(data = as.matrix(test_data %>% select(-Churn)), label = test_data$Churn)
\end{lstlisting}

\textbf{Explanation:}
\begin{itemize}
    \item \texttt{xgb.DMatrix()} -- Special data format that XGBoost needs.
    \item \texttt{data = as.matrix(...)} -- Convert features to matrix format.
    \item \texttt{label = ...} -- The target variable (what we're predicting).
\end{itemize}

\subsection{Lines 88-95: XGBoost Parameters}

\begin{lstlisting}[language=R]
  params <- list(
    objective = "binary:logistic",
    eta = 0.01,
    max_depth = 3,
    eval_metric = "auc",
    subsample = 0.8,
    colsample_bytree = 0.8
  )
\end{lstlisting}

\textbf{Explanation:}

\begin{table}[h]
\centering
\begin{tabular}{@{}lll@{}}
\toprule
\textbf{Parameter} & \textbf{Value} & \textbf{Meaning} \\ \midrule
objective & binary:logistic & Binary classification, output probability \\
eta & 0.01 & Learning rate (smaller = slower but more accurate) \\
max\_depth & 3 & Maximum tree depth (prevents overfitting) \\
eval\_metric & auc & Evaluate using Area Under ROC Curve \\
subsample & 0.8 & Use 80\% of data for each tree \\
colsample\_bytree & 0.8 & Use 80\% of features for each tree \\ \bottomrule
\end{tabular}
\end{table}

\subsection{Line 96: Training the Model}

\begin{lstlisting}[language=R]
  model <- xgb.train(params = params, data = dtrain, nrounds = 1000, verbose = 0)
\end{lstlisting}

\textbf{Explanation:}
\begin{itemize}
    \item \texttt{xgb.train()} -- Train XGBoost model.
    \item \texttt{params} -- The parameters we defined.
    \item \texttt{data = dtrain} -- Training data.
    \item \texttt{nrounds = 1000} -- Build 1000 trees (boosting rounds).
    \item \texttt{verbose = 0} -- Don't print progress messages.
\end{itemize}

\subsection{Lines 97-107: Finding Best Threshold}

\begin{lstlisting}[language=R]
  probs <- predict(model, dtest)
  thresholds <- seq(0.35, 0.65, by = 0.005)
  local_best_acc <- 0
  local_best_thresh <- 0.5
  for (t in thresholds) {
    preds <- ifelse(probs > t, 1, 0)
    acc <- mean(preds == test_data$Churn)
    if (acc > local_best_acc) {
      local_best_acc <- acc
      local_best_thresh <- t
    }
  }
\end{lstlisting}

\textbf{Explanation:}
\begin{itemize}
    \item \textbf{Line 97:} \texttt{predict()} -- Get probability predictions (0.0 to 1.0).
    \item \textbf{Line 98:} \texttt{seq(0.35, 0.65, by = 0.005)} -- Create sequence of thresholds to try.
    \item \textbf{Lines 101-107:} Try each threshold:
    \begin{itemize}
        \item If probability > threshold, predict 1 (churn), else 0 (no churn).
        \item Calculate accuracy for this threshold.
        \item Keep track of which threshold gives best accuracy.
    \end{itemize}
\end{itemize}

\textbf{Why different thresholds?} Default is 0.5, but sometimes 0.45 or 0.55 works better depending on the data.

\subsection{Lines 109-120: Tracking Best Overall Model}

\begin{lstlisting}[language=R]
  if (local_best_acc > best_overall_acc) {
    best_overall_acc <- local_best_acc
    best_thresh <- local_best_thresh
    best_model <- model
    best_probs <- probs
    best_actual <- test_data$Churn
    best_seed <- seed_val
  }
  if (best_overall_acc >= 0.86) {
    cat("\n TARGET done Stopping search.\n")
    break
  }
}
\end{lstlisting}

\textbf{Explanation:}
\begin{itemize}
    \item \textbf{Lines 109-116:} If this seed's accuracy beats our best so far, save all its results.
    \item \textbf{Lines 117-120:} If we achieved 86\% accuracy, stop searching (we met our goal).
    \item \texttt{break} -- Exit the for loop early.
\end{itemize}

% ============================================================
\section{Final Evaluation}
% ============================================================

\subsection{Lines 122-127: Calculating Metrics}

\begin{lstlisting}[language=R]
cat("\n>>> Section 3.1: Beyond Accuracy: A Deep Dive into the Confusion Matrix...\n")
final_preds <- ifelse(best_probs > best_thresh, 1, 0)
cm <- confusionMatrix(as.factor(final_preds), as.factor(best_actual), positive = "1")
precision <- posPredValue(as.factor(final_preds), as.factor(best_actual), positive = "1")
recall <- sensitivity(as.factor(final_preds), as.factor(best_actual), positive = "1")
f1 <- 2 * ((precision * recall) / (precision + recall))
\end{lstlisting}

\textbf{Explanation:}
\begin{itemize}
    \item \textbf{Line 123:} Make final predictions using best threshold.
    \item \textbf{Line 124:} \texttt{confusionMatrix()} -- Creates confusion matrix comparing predictions vs actual.
    \item \textbf{Line 125:} \texttt{posPredValue()} = Precision = Of all predicted churners, how many actually churned?
    \item \textbf{Line 126:} \texttt{sensitivity()} = Recall = Of all actual churners, how many did we catch?
    \item \textbf{Line 127:} F1 Score = Harmonic mean of precision and recall (balance between both).
\end{itemize}

The F1 Score formula is:
\[
F1 = 2 \times \frac{Precision \times Recall}{Precision + Recall}
\]

\subsection{Lines 128-136: Printing Results}

\begin{lstlisting}[language=R]
cat(rep("-", 50), "\n")
cat(sprintf("WINNING SEED:        %d\n", best_seed))
cat(sprintf("FINAL ACCURACY:      %.2f%%\n", best_overall_acc * 100))
cat(sprintf("PRECISION:           %.4f\n", precision))
cat(sprintf("RECALL:              %.4f\n", recall))
cat(sprintf("F1 SCORE:            %.4f\n", f1))
cat(rep("-", 50), "\n")
cat("\nConfusion Matrix Table:\n")
print(cm$table)
\end{lstlisting}

\textbf{Explanation:}
\begin{itemize}
    \item Prints all the evaluation metrics in a formatted way.
    \item \texttt{cm\$table} -- The actual confusion matrix table.
\end{itemize}

% ============================================================
\section{Summary of the Workflow}
% ============================================================

\begin{enumerate}
    \item \textbf{INSTALL \& LOAD PACKAGES} -- Get all required tools ready
    \item \textbf{LOAD DATA} -- Read the Telco customer data CSV file
    \item \textbf{CLEAN DATA}
    \begin{itemize}
        \item Fix missing values in TotalCharges
        \item Simplify categorical values
        \item Remove customerID column
        \item Convert all text to numbers
    \end{itemize}
    \item \textbf{EXPLORE DATA (EDA)}
    \begin{itemize}
        \item Calculate churn rate
        \item Create 4 visualizations
    \end{itemize}
    \item \textbf{BALANCE DATA (SMOTE)} -- Create synthetic samples to balance churners/non-churners
    \item \textbf{TRAIN MODEL}
    \begin{itemize}
        \item Try multiple random seeds
        \item For each seed: split data, train XGBoost, find best threshold
        \item Keep the model that achieves 86\%+ accuracy
    \end{itemize}
    \item \textbf{EVALUATE}
    \begin{itemize}
        \item Calculate precision, recall, F1 score
        \item Show confusion matrix
    \end{itemize}
\end{enumerate}

% ============================================================
\section{Key Concepts Explained}
% ============================================================

\subsection{What is Churn?}
Churn = When a customer leaves/cancels their service. We want to predict which customers might churn so the company can try to keep them.

\subsection{What is XGBoost?}
XGBoost (eXtreme Gradient Boosting) is a powerful machine learning algorithm that builds many decision trees, where each new tree tries to fix the mistakes of the previous trees.

\subsection{What is SMOTE?}
SMOTE creates synthetic (fake but realistic) examples of the minority class by interpolating between existing examples. This helps the model learn to recognize churners better.

\subsection{What is a Confusion Matrix?}
A table showing prediction results:

\begin{table}[h]
\centering
\begin{tabular}{@{}lcc@{}}
\toprule
 & \textbf{Predicted No} & \textbf{Predicted Yes} \\ \midrule
\textbf{Actual No} & True Negative (TN) & False Positive (FP) \\
\textbf{Actual Yes} & False Negative (FN) & True Positive (TP) \\ \bottomrule
\end{tabular}
\end{table}

\subsection{Evaluation Metrics}

\begin{itemize}
    \item \textbf{Accuracy} = $\frac{TP + TN}{Total}$ = Overall correctness
    \item \textbf{Precision} = $\frac{TP}{TP + FP}$ = When we predict churn, how often are we right?
    \item \textbf{Recall} = $\frac{TP}{TP + FN}$ = Of all churners, how many did we catch?
    \item \textbf{F1 Score} = $2 \times \frac{Precision \times Recall}{Precision + Recall}$ = Balance between Precision and Recall
\end{itemize}

\end{document}
